

\section{Longest Palindromic Subsequence (with long-form proof)}\label{dp: sec: LPS}\doHeader{Longest Palindrome}

\begin{mylist}
\item[\bf input:]
  A string $X[1..n]$ (a sequence of characters).
  
\item[\bf output:]
  The maximum length of any subsequence of $X$ that is a palindrome.
  (A palindrome is a sequence that is the same as its reverse.
  The subsequence need not be contiguous.
  For example, ``yobananaboy'' is a palindromic subsequence
  of ``your urban cabana bat toy''.)

\end{mylist}


\noindent
Here is the algorithm.  Fix any input instance $X[1..n]$.

\paragraph{subproblems and recurrence.}
% [review 2026-09-27] fix: M(i,j) was defined only for 1 <= i <= j <= n, but the recurrence uses M(i+1,i) (empty substrings); extended the domain to j >= i-1, with M(i,i-1) = 0
For each $i$ and $j$ with $1\le i \le n+1$ and $i-1\le j \le n$,
define $M(i, j)$ to be the maximum length of any palindromic subsequence of $X[i..j]$.
(When $j = i-1$, $X[i..j]$ is the empty string, so $M(i, i-1) = 0$.)
The given instance $X$ has solution $M(1, n)$.
The recurrence relation is
\[M(i, j)  =
\begin{cases}
  0 & \text{if } j < i, \\
  1 & \text{if } j = i, \\
  \max\big(M(i, j-1), M(i+1, j), 2 + M(i+1, j-1)\big) & \text{if } j>i \text{ and } X[i] = X[j], \\
  \max\big(M(i, j-1), M(i+1, j)\big) & \text{if } j>i \text{ and } X[i] \ne X[j].
\end{cases}
\]
(We can simplify this somewhat per the comment at the end below.)

\paragraph{correctness.}
The base cases hold by inspection.  Here's the intuition for the remaining cases.
Consider the palindromic subsequences of $X[i..j]$.
Categorize these into three types:
(A) those that include $X[i]$ and $X[j]$ (necessarily starting with $X[i]$ and ending with $X[j]$),
(B) those that don't include $X[i]$,
(C) those that don't include $X[j]$.
Every subsequence is of one (or more) of these types.
So the maximum length of any such subsequence
is the maximum length of any type-A, type-B, or type-C subsequence.

First consider the case that $X[i] = X[j]$.
Then those of type A are of the form $X[i], S, X[j]$
where $S$ is any palindromic subsequence of $X[i+1..j-1]$;
so the longest of type A has length $2+M(i+1, j-1)$.  
% [review 2026-09-27] fix: types B and C were swapped relative to their definitions above (B = doesn't include X[i], C = doesn't include X[j]); swapped the substrings and M-terms
Those of type B are just palindromic subsequences of $X[i+1..j]$;
the longest of these has length $M(i+1, j)$.
Those of type C are just palindromic subsequences of $X[i..j-1]$;
the longest of these has length $M(i, j-1)$.
So the recurrence holds in this case.

In the remaining case ($X[i] \ne X[j]$),
there are none of type A, because any palindromic subsequence
has to start and end with the same character.
The analysis of the other two types is the same as in the case $X[i] = X[j]$.
So the recurrence holds in this case.
A long-form proof of correctness is on the next page.

\paragraph{time.}
There are $O(n^2)$ subproblems.
For each, the right-hand side of the recurrence can be evaluated in constant time.
So the algorithm (iterative, or memoized recursive) takes time $\Theta(n^2)$.

\begin{theorem}
  % [review 2026-09-27] fix: theorem said O(nm) but there is only one string of length n (time analysis says Theta(n^2)); changed to O(n^2)
  There is an $O(n^2)$-time algorithm for \prob{Longest Palindromic Subsequence}.
\end{theorem}

\paragraph{reduction to Longest Path.}
Consider the DAG with a node for each subproblem
and a directed edge for each dependency (with appropriate weight 0 or 2),
% [review 2026-09-27] fix: base-case edges were wrong (M(i,i+1) is not a base case; M(i,i)=1 needs a weight-1 edge; the empty M(i+1,i)=0 need weight-0 edges) and the target M(n,m) should be M(1,n)
and additional directed edges from the node for $M(0,0)$
to the nodes for each $M(i, i)$ with $1\le i\le n$ (of weight 1)
and to the nodes for each $M(i+1, i)$ with $0\le i\le n$ (the empty substrings, of weight 0).
The problem is equivalent to finding a longest path in this DAG from the node for $M(0,0)$
to the node for $M(1, n)$.

\paragraph{simplifying the recurrence.}
In the case that $X[i]=X[j]$,
for any type-B or type-C subsequence,
there is always a type-A subsequence that is at least as long.
(We can use a ``greedy'' exchange argument to prove this; we omit the details.)
So the maximum in the recurrence for this case is always
achieved by the third term, $2+M(i+1, j-1)$,
and the recurrence simplifies to
\[
  M(i, j)  =
  \begin{cases}
    0 & \text{if } j < i, \\
    1 & \text{if } j = i, \\
    2 + M(i+1, j-1) & \text{if } j>i \text{ and } X[i] = X[j], \\
    \max\big(M(i, j-1), M(i+1, j)\big) & \text{if } j>i \text{ and } X[i] \ne X[j].
  \end{cases}
\]

\subsection*{Long-form proof of correctness}
Here is a long-form proof that the first recurrence is correct.

\begin{longFormProof}%[First long-form proof.]
  \step Consider any $i$ and $j$ with $1\le i\le j\le n$.  Note that $j-i+1$ is the length of the string $X[i..j]$.

  \smallskip 
  \begin{block}
    \step \underline{Case 1.} In the case that $j-i+1\le 1$, the string $X[i..j]$ has length one or zero.
    \step So it is it's own longest palindromic subsequence, $M(i, j)=j-i+1$, and the recurrence holds.
  \end{block}

  \smallskip 
  \begin{block}
    \step \underline{Case 2.} Next consider the case that $X[i]=X[j]$ (and $j-i+1\ge 2$).
    \step Consider all possible palindromic subsequences of $X[i..j]$.
    \step Categorize these into three types:
    \begin{mylist}
    \item[(A)] those that include $X[i]$ and $X[j]$,
    \item[(B)] those that don't include $X[i]$, and
    \item[(C)] those that don't include $X[j]$.
    \end{mylist}

    \step Each palindromic subsequence of $X[i..j]$ is of  one (or more!) of these types.

    \step\label{dp: step: A or B or C} So $M(i, j)$ is the maximum length of any type-A, type-B, or type-C subsequence.

    \smallskip 
    
    \step The type-A's are of the form $X[i], s, X[j]$,
    where $s$ is any palindromic subseq.~of $X[i+1..j-1]$.

    \step The length of such a subsequence $X[i], s, X[j]$ is maximized when the length of $s$ is maximized.

    \step By definition of $M(i+1,j-1)$, the maximum length of any such $s$ is $M(i+1, j-1)$.

    \step\label{dp: step: palindrome A}
    So the maximum length of any type-A subsequence is $2+M(i+1, j-1)$.

    \smallskip
    
    \step The type-B subsequences are just palindromic subsequences of $X[i+1..j]$.

    % [review 2026-09-27] fix: cited "definition of M(i-1,j)" but the quantity is M(i+1,j); changed i-1 to i+1
    \step\label{dp: step: palindrome B}
    So (by definition of $M(i+1,j)$) the maximum length of any type-B subsequence is $M(i+1, j)$.

    \smallskip
    
    \step\label{dp: step: C}
    Similarly, the maximum length of any type-C subsequence is $M(i, j-1)$.
    
    \smallskip 

    \step By Steps~\ref{dp: step: A or B or C},~\ref{dp: step: palindrome A},~\ref{dp: step: palindrome B}, and~\ref{dp: step: C},
    $M(i, j) = \max\big(2+M(i+1, j-1), M(i+1, j), M(i, j-1)\big)$.

    \step So the recurrence holds in this case.
  \end{block}

  \smallskip 
  \begin{block}
    % [review 2026-09-27] fix: length condition was written i-j+1 >= 2 (never true for i <= j); changed to j-i+1 >= 2 as in Case 2
    \step \underline{Case 3.}  Finally consider the remaining case, when $X[i] \ne X[j]$ and $j-i+1\ge 2$.
    \step Since $X[i]\ne X[j]$, no palindromic subsequence of $X[i..j]$ starts with $X[i]$ and ends with $X[j]$.
    \step So, the reasoning in Case 2 applies, except that there are no type-A subsequences.
    \step And by that reasoning, $M(i, j) = \max\big(M(i+1, j), M(i, j-1)\big)$.
    \step So the recurrence holds in this case.
  \end{block}

  \smallskip
  \step By the case analysis (Cases 1--3) above, the recurrence holds.
\end{longFormProof}

\pythonCode{dynamic_programming/longest_palindrome.py}


% \vfill\emph{A more detailed long-form proof, using a different approach, is on the next page.}\newpage

% \begin{longFormProof}[Second long-form proof.]
%   \step Consider any $i$ and $j$ with $1\le i\le j\le n$.  Note that $j-i+1$ is the length of the string $X[i..j]$.

%   \smallskip 
%   \begin{block}
%     \step \underline{Case 1.} In the case that $j-i+1\le 1$, the string $X[i..j]$ has length one or zero.
%     \step So it is it's own longest palindromic subsequence, $M(i, j)=j-i+1$, and the recurrence holds.
%   \end{block}

%   \smallskip 
%   \begin{block}
%     \step \underline{Case 2.} Next consider the case that $j-i+1\ge 2$.

%     \medskip

%     \step First we show that $M(i, j)$ is \emph{at least} the right-hand side of the recurrence (the RHS).

%     \step Let $s$ be a palindromic subsequence of $X[i+1, j]$ with length $|s|$ equal to $M(i+1, j)$.
%     \\ (Such an $s$ exists by definition of $M(i+1, j-1)$.)

%     \step Then $s$ is also a palindromic subsequence of $X[i, j]$.
    
%     \step $M(i, j)$ is the maximum length of any palindromic subsequence of $X[i..j]$.

%     \step\label{dp: step: 1} So $M(i, j)$ is at least the length of $s$.  That is, $M(i, j) \ge |s| = M(i+1, j)$.
    
%     \smallskip

%     \step\label{dp: step: 2} By similar reasoning, $M(i, j) \ge M(i, j-1)$.

%     \smallskip 

%     \begin{block}
%       \step\underline{Case 2.1}
%       Now consider the case that $X[i] \ne X[j]$.

%       \step\label{dp: step: 12} By Steps~\ref{dp: step: 1} and~\ref{dp: step: 2}, $M(i, j) \ge \max\big(M(i+1, j), M(i, j-1)\big)$.

%       \step That is (by inspection of the recurrence, using $X[i] \ne X[j]$), $M(i, j)$ is at least the RHS.
%     \end{block}
    
%     \smallskip 

%     \begin{block}
%       \step\underline{Case 2.2}
%       In the remaining case $X[i] = X[j]$.

%       \step Let $s$ be a palindromic subsequence of $X[i+1..j-1]$ with length $|s|$ equal to $M(i+1, j-1)$.
%       \\ (Such an $s$ exists by definition of $M(i+1, j-1)$.)

%       \step Then $X[i], s, X[j]$ is a palindromic subsequence of $X[i..,j]$.

%       \step\label{dp: step: 3} So $M(i, j)$ is at least the length of $X[i], s, X[j]$.
%       That is, $M(i, j) \ge 2+|s| = 2 + M(i+1, j-1)$.

%       \smallskip 

%       \step\label{dp: step: 123}
%       By Steps~\ref{dp: step: 1},~\ref{dp: step: 2}, and~\ref{dp: step: 3},
%       $M(i, j) \ge \max\big(M(i+1, j), M(i, j-1), 2 + M(i+1, j-1)\big)$.

%       \step That is (by inspection of the recurrence, using $X[i]=X[j]$), $M(i, j)$ is at least the RHS.
%     \end{block}
    
%     \smallskip 

%     \step\label{dp: step: 4} By the case analysis in Cases 2.1 and 2.2, $M(i, j)$ is at least the RHS.

%     \smallskip 

%     \step Next we show that $M(i, j)$ is \emph{at most} the right-hand side of the recurrence.

%     \step Let $s$ be a palindromic subsequence of $X[i..j]$ of length $M(i, j)$.
%     \\(Such an $s$ exists by definition of $M(i, j)$.)

%     \begin{block}
%       \step \underline{Case 2.A}
%       First consider the case that $s$ does not start with $X[i]$.

%       \step Then $s$ is a palindromic subsequence of $X[i+1..j]$.

%       \step So $s$ has length at most $M(i+1, j)$.  That is $M(i, j) = |s| \le M(i+1, j)$.

%       \step So, by inspection of the recurrence, $M(i, j)$ is at most the RHS.
%     \end{block}

%     \smallskip 
    
%     \begin{block}
%       \step \underline{Case 2.B}
%       Next consider the case that $s$ does not end with $X[j]$.

%       \step By the same reasoning as in Case 2.A, $M(i, j) = |s| \le M(i,j-1)$.
      
%       \step Again, by inspection of the recurrence, $M(i, j)$ is at most the RHS.
%     \end{block}

%     \smallskip 
    
%     \begin{block}
%       \step \underline{Case 2.C}
%       In the remaining case $s$ starts with $X[i]$ and ends with $X[j]$

%       \step So $X[i]=X[j]$.

%       \step Also, $s$ is of the form $X[i], s', X[j]$ where $s'$
%       is a palindromic subsequence of $X[i+1..j-1]$.

%       \step So $|s'| \le M(i+1, j-1)$.  It follows that $M(i, j) = |s| = |s'| + 2 \le M(i+1, j-1) + 2$.

%       \step So, by inspection of the recurrence (using $X[i] \ne X[j]$), $M(i, j)$ is at most the RHS.
%     \end{block}

%     \smallskip 

%     \step By the case analysis (Cases 2.A-C) above,  $M(i, j)$ is at most the RHS. 

%     \step By this and Step~\ref{dp: step: 4},  $M(i, j)$ equals the RHS (in Case 2).
%   \end{block}

%   \step By the case analysis (Cases 1 and 2) above,  $M(i, j)$ equals the RHS.
% \end{longFormProof}

% \noindent
% This algorithm is equivalent to running the \prob{Longest Path} algorithm
% on an appropriate DAG.

% \paragraph{External resources on \prob{Knapsack}}

% \begin{itemize}
% % \item CLRS Chapter 15. 

% \item CLRS Chapter 16, Exercise 16.2--2 

% \item Dasgupta et al.~Chapter 6.

% \item Kleinberg \& Tardos Chapter 6.

% % \item Erickson's Lecture 5
% %
% %     \url{https://jeffe.cs.illinois.edu/teaching/algorithms/notes/05-dynprog.pdf}
  
% % \item MIT lecture videos 

% \item MIT lecture video
  
% \newcommand{\URL}[1]{\par-- \parbox[t]{0.8\textwidth}{\url{#1}} \medskip\par}


% \URL{https://ocw.mit.edu/courses/electrical-engineering-and-computer-science/6-006-introduction-to-algorithms-fall-2011/lecture-videos/lecture-19-dynamic-programming-i-fibonacci-shortest-paths/}


% \URL{https://ocw.mit.edu/courses/electrical-engineering-and-computer-science/6-006-introduction-to-algorithms-fall-2011/lecture-videos/lecture-20-dynamic-programming-ii-text-justification-blackjack}

    % \URL{https://ocw.mit.edu/courses/electrical-engineering-and-computer-science/6-006-introduction-to-algorithms-fall-2011/lecture-videos/lecture-21-dp-iii-parenthesization-edit-distance-knapsack}

%   \URL{https://ocw.mit.edu/courses/electrical-engineering-and-computer-science/6-006-introduction-to-algorithms-fall-2011/lecture-videos/lecture-22-dp-iv-guitar-fingering-tetris-super-mario-bros}
  
% \item \prob{Seam Carving}:
%   \url{https://en.wikipedia.org/wiki/Seam_carving},
%   CLRS Problem 15-8.

% \end{itemize}

%%% Local Variables:
%%% mode: latex
%%% TeX-master: "dynamic_programming"
%%% End:
